% TOP_PROBLEM: {"id": 1469, "title": "Serre's Positivity Conjecture", "category_id": 4, "category": {"id": 4, "name": "algebra", "display_name": "Algebra", "description": "Group theory, ring theory, field theory, and algebraic structures.", "slug": "algebra", "order_index": 4, "created_at": "2026-07-31T15:26:25.671Z"}, "status": "open"}
% FIRST_POSTED: 2026-09-25
% ENTRY_KIND: statement_only
% !TeX program = lualatex
% Definition and sources only; this file is not an LLM solution attempt.
\documentclass[11pt]{article}
\usepackage[margin=25mm]{geometry}
\usepackage{fontspec}
\setmainfont{Latin Modern Roman}
\usepackage{amsmath,amssymb,mathtools}
\usepackage{unicode-math}
\setmathfont{Latin Modern Math}
\usepackage{xurl,hyperref}
\hypersetup{colorlinks=true,urlcolor=blue}
\setcounter{secnumdepth}{0}
\setlength{\parindent}{0pt}
\setlength{\parskip}{5pt}
\setlength{\emergencystretch}{3em}
\begin{document}
\section*{\texorpdfstring{Serre's Positivity Conjecture}{Serre's Positivity Conjecture}}\label{1469}
\subsection{Definitions and mathematical statement}
Let $(R,\mathfrak m)$ be a commutative regular Noetherian local ring and let $M,N$ be nonzero finitely generated modules with $M\otimes_RN$ of finite length. Their dimensions are dimensions of their supports. Since $R$ has finite global dimension, the alternating sum
\[
 \chi^R(M,N)=\sum_{i\ge0}(-1)^i
 \operatorname{length}_R\operatorname{Tor}_i^R(M,N)
\]
is finite; the support hypothesis makes each term finite. Serre positivity asserts $\chi^R(M,N)>0$ whenever $\dim M+\dim N=\dim R$. No equal-characteristic or unramified assumption is added. Strict positivity is stronger than nonnegativity, while the complementary dimension regimes are not part of this selected statement.
\par\smallskip\textit{Formulation and concept references:} \href{https://arxiv.org/html/2404.17572}{[S1]}.

\subsection{Short English statement}
When two modules over a regular local ring meet in finite length and their dimensions add correctly, must their intersection multiplicity be strictly positive?
\subsection{Sources}
\begin{itemize}
\item[S1]N. KC and A. J. Soto Levins, On liftings of modules of finite projective dimension, Section 1.1 (2024).
\url{https://arxiv.org/html/2404.17572}.
\textit{Formulation source.}
\end{itemize}
\end{document}
